% \iffalse meta-comment
% (The MIT License)
%
% Copyright (c) 2021-2022 Yegor Bugayenko
%
% Permission is hereby granted, free of charge, to any person obtaining a copy
% of this software and associated documentation files (the 'Software'), to deal
% in the Software without restriction, including without limitation the rights
% to use, copy, modify, merge, publish, distribute, sublicense, and/or sell
% copies of the Software, and to permit persons to whom the Software is
% furnished to do so, subject to the following conditions:
%
% The above copyright notice and this permission notice shall be included in all
% copies or substantial portions of the Software.
%
% THE SOFTWARE IS PROVIDED 'AS IS', WITHOUT WARRANTY OF ANY KIND, EXPRESS OR
% IMPLIED, INCLUDING BUT NOT LIMITED TO THE WARRANTIES OF MERCHANTABILITY,
% FITNESS FOR A PARTICULAR PURPOSE AND NONINFRINGEMENT. IN NO EVENT SHALL THE
% AUTHORS OR COPYRIGHT HOLDERS BE LIABLE FOR ANY CLAIM, DAMAGES OR OTHER
% LIABILITY, WHETHER IN AN ACTION OF CONTRACT, TORT OR OTHERWISE, ARISING FROM,
% OUT OF OR IN CONNECTION WITH THE SOFTWARE OR THE USE OR OTHER DEALINGS IN THE
% SOFTWARE.
% \fi

%%% \CheckSum{0}
%
% \CharacterTable
%  {Upper-case    \A\B\C\D\E\F\G\H\I\J\K\L\M\N\O\P\Q\R\S\T\U\V\W\X\Y\Z
%   Lower-case    \a\b\c\d\e\f\g\h\i\j\k\l\m\n\o\p\q\r\s\t\u\v\w\x\y\z
%   Digits        \0\1\2\3\4\5\6\7\8\9
%   Exclamation   \!     Double quote  \"     Hash (number) \#
%   Dollar        \$     Percent       \%     Ampersand     \&
%   Acute accent  \'     Left paren    \(     Right paren   \)
%   Asterisk      \*     Plus          \+     Comma         \,
%   Minus         \-     Point         \.     Solidus       \/
%   Colon         \:     Semicolon     \;     Less than     \<
%   Equals        \=     Greater than  \>     Question mark \?
%   Commercial at \@     Left bracket  \[     Backslash     \\
%   Right bracket \]     Circumflex    \^     Underscore    \_
%   Grave accent  \`     Left brace    \{     Vertical bar  \|
%   Right brace   \}     Tilde         \~}

% \GetFileInfo{huawei.dtx}
% \DoNotIndex{\endgroup,\begingroup,\let,\else,\fi,\newcommand,\newenvironment}

% \iffalse
%<*driver>
\ProvidesFile{huawei.dtx}
%</driver>
%<class>\NeedsTeXFormat{LaTeX2e}
%<class>\ProvidesClass{huawei}
%<*class>
[2022-10-14 0.14.2 Template for Huawei Documents]
%</class>
%<*driver>
\documentclass{ltxdoc}
\usepackage[tt=false, type1=true]{libertine}
\usepackage[T1]{fontenc}
\usepackage[utf8]{inputenc}
\usepackage{href-ul}
\PageIndex
\EnableCrossrefs
\CodelineIndex
\RecordChanges
\begin{document}
	\DocInput{huawei.dtx}
	\PrintChanges
	\PrintIndex
\end{document}
%</driver>
% \fi

% \title{\LaTeX{} Class |huawei|\thanks{The sources are in GitHub at \href{https://github.com/yegor256/huawei.cls}{yegor256/huawei.cls}}}
% \author{Yegor Bugayenko \\ \texttt{yegor256@gmail.com}}
% \date{\filedate, \fileversion}
%
% \maketitle
%
% \section{Introduction}
%
% The provided class |huawei| helps you design your work
% documents and presentations keeping the code short and the
% style elegant enough both for management and technical
% papers. To use the class you simply mention its name
% in the preamble:
%
%\iffalse
%<*verb>
%\fi
\begin{verbatim}
\documentclass{huawei}
\begin{document}
Hello, world!
\end{document}
\end{verbatim}
%\iffalse
%</verb>
%\fi

% We recommend you to use
% |latexmk| to compile your |.tex| files to |.pdf|.
% The simplest setup will require a few files staying next to your
% |.tex| file, in the same directory (``story'' is the name
% of your project here):

%\iffalse
%<*verb>
%\fi
\begin{verbatim}
story\
  .latexmkrc
  .gitignore
  story.tex
\end{verbatim}
%\iffalse
%</verb>
%\fi

% The content of your |.latexmkrc| file would be this:

%\iffalse
%<*verb>
%\fi
\begin{verbatim}
$pdflatex = 'pdflatex %O --shell-escape %S';
\end{verbatim}
%\iffalse
%</verb>
%\fi

% Make sure the |.gitignore| file lists all the files generated
% by |pdflatex| during the compilation. You don't need to commit
% them to your repository, since they are temporary and will
% be generated again when you compile your document.

% In order to compile the document, just say |latexmk| on the command line.

% \section{Class Options}

% There are a few class options, provided in square brackets after the
% |\documentclass|, which can help you fine-tune
% the layout of your document:

% \DescribeMacro{landscape}
%   makes the document in landscape format, also changing the size
%   of the paper to 16x9 inches (the default page size is
%   \href{https://en.wikipedia.org/wiki/Paper_size}{A4}),
%   making it perfect for presentations.

% \DescribeMacro{dark}
% Turns on a dark layout, where the page color is black and the text
% is white.

% \DescribeMacro{slides}
%   makes all headers a bit larger, assuming that the document
%   is in the landscape mode and being presented as a slide deck.

% \DescribeMacro{nocover}
%   avoid printing the cover images on the first page by the
%   |\PrintTitlePage| command.

% \DescribeMacro{anonymous}
%   removes the name of the author everywhere, including the bottom
%   of the page, where the author's name stays next to the name of the
%   company.

% \DescribeMacro{nobrand}
%   avoid mentioning the brand of Huawei anywhere
%   in the document and removes the logo too.

% \DescribeMacro{nosecurity}
%   avoids mentioning the level of security at the right top
%   corner of the document and also avoids showing the ID of the author
%   where it usually is visible.

% \DescribeMacro{nodate}
%   don't show the date and time at the bottom of each page,
%   where they usually are rendered in ISO~8601 format.

% \DescribeMacro{nopaging}
%   avoids page numbers at the bottom of each page.

% \DescribeMacro{authordraft}
%   prints a big ``It's a draft'' message across each page.

% \section{Preamble}

% In the preamble you can specify meta information about the document,
% such as its title or author's name, here is how:

%\iffalse
%<*verb>
%\fi
\begin{verbatim}
\documentclass{huawei}
\renewcommand*\thetitle{Making Compression 15\% Faster}
\renewcommand*\thesubtitle{Technical Report}
\renewcommand*\theauthor{Yegor Bugayenko}
\begin{document}
\maketitle
Hello, world!
\end{document}
\end{verbatim}
%\iffalse
%</verb>
%\fi

% It's recommended to use |\renewcommand*| instead of
% |\renewcommand| in order to let \LaTeX{} catch you
% if by mistake a new line gets into the content.

% The following meta commands are defined:

% \DescribeMacro{\thetitle}
% |\thetitle| is the main title of the document
%   to be used in the text and in the properties of the PDF document.

% \DescribeMacro{\thesubtitle}
% |\thesubtitle| is the subtitle to be
%   printed under the title.

% \DescribeMacro{\theauthor}
% |\theauthor| is the author of the document
%   in ``first-name last-name'' format.

% \DescribeMacro{\theid}
% |\theid| is the internal ID of the author, if
%   it's applicable.

% \DescribeMacro{\thesecurity}
% It is the level of security of
%   the document, which is usually printed at the top right
%   corner of it; usual values are ``Internal,''
%   ''Confidential,'' or ``Secret.''

% Default values of all these commands are empty. If you don't
% renew them in your document, nothing will be printed.

% \section{Custom Commands}

% Inside the document body you can use these commands:

% \DescribeMacro{\PrintFirstPage}
% |\PrintFirstPage|
%   prints the first page of a project charter or a similar landscape documents,
%   placing the image |front-image.pdf| on the front (the file should be present
%   in the current dir. If you don't have the front image file, just leave
%   the first argument empty.

% \DescribeMacro{\PrintLastPage}
% |\PrintLastPage|
%   prints the last page of a project charter or a similar landscape document.

% \DescribeMacro{\PrintThankYouPage}
% |\PrintThankYouPage|
%   prints the last page with a "Thank You" message in the center.

% \DescribeMacro{\PrintDisclaimer}
% |\PrintDisclaimer|
%   prints a paragraph at the bottom of the page with a standard disclaimer.

% \section{Best Practices}

% You are free to design your documents any way you want. However,
% it would be convenient for yourself and for your readers, if you follow
% the convention we have for business and technical documents.

% The rule of thumb is simple: try \emph{not} to format your documents.
% Instead, let the class designed by us do this work for you. Just type
% the content without changing the layout, adding colors, changing fonts,
% etc. The less you modify the look-and-feel, the better your documents
% will be perceived by your readers.

% \subsection{Two Columns}

% In the landscape format it's recommended to use two columns, for better
% readability of the text. Here is how:

%\iffalse
%<*verb>
%\fi
\begin{verbatim}
\documentclass{huawei}
\begin{document}
\newpage
\begin{multicols}{2}
\section*{First}
Here goes the first column content.
\columnbreak
\section*{Second}
Here goes the second column content.
\end{multicols}
\end{document}
\end{verbatim}
%\iffalse
%</verb>
%\fi

% \subsection{Crumbs}

% When you need to put many small information pieces into one page,
% we recommend you to use ``crumbs'':

%\iffalse
%<*verb>
%\fi
\begin{verbatim}
\documentclass{huawei}
\begin{document}
\newpage
\section*{Project Details}
\begin{multicols}{2}
\raggedright
\PrintCrumb{Budget}{\$100K}

\PrintCrumb{Duration}{5 months}
\end{multicols}
\end{document}
\end{verbatim}
%\iffalse
%</verb>
%\fi

% \StopEventually{}

% \section{Implementation}

% \changes{v0.1.0}{2021/04/18}{Initial version}

% First, we load the parent class:
%    \begin{macrocode}
\LoadClass[12pt]{article}
%    \end{macrocode}

% Then, we parse package options:
%    \begin{macrocode}
\makeatletter
\newif\ifhuawei@dark
\DeclareOption{dark}{\huawei@darktrue}
\newif\ifhuawei@slides
\DeclareOption{slides}{\huawei@slidestrue}
\newif\ifhuawei@nosecurity
\DeclareOption{nosecurity}{\huawei@nosecuritytrue}
\newif\ifhuawei@authordraft
\DeclareOption{authordraft}{\huawei@authordrafttrue}
\newif\ifhuawei@nobrand
\DeclareOption{nobrand}{\huawei@nobrandtrue}
\newif\ifhuawei@nodate
\DeclareOption{nodate}{\huawei@nodatetrue}
\newif\ifhuawei@nocover
\DeclareOption{nocover}{\huawei@nocovertrue}
\newif\ifhuawei@nopaging
\DeclareOption{nopaging}{\huawei@nopagingtrue}
\newif\ifhuawei@landscape
\DeclareOption{landscape}{\huawei@landscapetrue}
\newif\ifhuawei@anonymous
\DeclareOption{anonymous}{\huawei@anonymoustrue}
\ProcessOptions\relax
\makeatother
%    \end{macrocode}

% Then, we set the size of the page:
%    \begin{macrocode}
\RequirePackage{geometry}
\makeatletter\ifhuawei@landscape
  \geometry{paperwidth=16in, paperheight=9in,
    left=4in, right=2in, top=1.5in, bottom=1.5in}
\else
  \geometry{a4paper, left=1.5in, right=1in,
    top=1.2in, bottom=1.2in}
\fi\makeatother
%    \end{macrocode}

% Then, we set the size of the font:
%    \begin{macrocode}
\makeatletter\ifhuawei@slides
  \usepackage[fontsize=24pt]{fontsize}
\fi\makeatother
%    \end{macrocode}

% Then, we include a important packages:
%    \begin{macrocode}
\RequirePackage[T1]{fontenc}
\RequirePackage[utf8]{inputenc}
\RequirePackage{tabularx}
\RequirePackage{anyfontsize}
\RequirePackage{multicol}
\RequirePackage{ragged2e}
\RequirePackage{multicol}
\RequirePackage[mmddyyyy,iso]{datetime}
  \newtimeformat{daytime}{\twodigit{\THEHOUR}:\twodigit{\THEMINUTE}}
\RequirePackage[tt=false,type1=true]{libertine}
\RequirePackage{paralist}
\RequirePackage{makecell}
\RequirePackage[para]{footmisc}
  \setlength{\footnotemargin}{2pt}
  \setlength{\footnotesep}{2pt}
\RequirePackage{graphicx}
\RequirePackage{array}
\RequirePackage[hidelinks]{hyperref} % to enable \pageref* command
\RequirePackage[abspath]{currfile}
\RequirePackage{enumitem}
  \setlist{nosep}
\RequirePackage{microtype}
\RequirePackage{wrapfig}
\RequirePackage{lastpage}
\RequirePackage{tikz}
\RequirePackage{tabularx}
\RequirePackage{titling}
\RequirePackage{svg}
\RequirePackage[absolute]{textpos}
  \TPGrid{16}{16}
%    \end{macrocode}

% Then, we configure |biblatex|:
%    \begin{macrocode}
\RequirePackage[maxnames=1,minnames=1,natbib=true,
  citestyle=numeric,bibstyle=numeric,doi=false,
  url=false,isbn=false,isbn=false]{biblatex}
%    \end{macrocode}

% Then, we set spacing between lines:
%    \begin{macrocode}
\RequirePackage{setspace}
\setstretch{1.08}
%    \end{macrocode}

% Then, we define branded colors:
%    \begin{macrocode}
\PassOptionsToPackage{table}{xcolor}
\RequirePackage{xcolor}
\definecolor{red}{HTML}{CF0A2C}
\definecolor{black}{HTML}{232527}
\definecolor{gray}{HTML}{878C8F}
\definecolor{yellow}{HTML}{F2DC5D}
\definecolor{blue}{HTML}{2274A5}
\definecolor{green}{HTML}{499167}
\definecolor{orange}{HTML}{F06543}
%    \end{macrocode}

% \begin{macro}{\huawei@header}
% Then, we define |\huawei@header|:
%    \begin{macrocode}
\makeatletter\newcommand\huawei@header[1]{{%
  \ifhuawei@slides%
    \setstretch{0.8}%
    \fontsize{19pt}{24pt}\selectfont%
  \else%
    \setstretch{0.8}%
    \fontsize{11pt}{13pt}\selectfont%
  \fi%
  \sffamily\color{gray}#1\par%
}}\makeatother
%    \end{macrocode}
% \end{macro}

% \begin{macro}{\huawei@bar}
% Then, we define |\huawei@bar|:
%    \begin{macrocode}
\makeatletter\newcommand\huawei@bar{%
  \begin{textblock}{1}[0,0](0,0)%
    \tikz \node[fill=red,minimum width=\TPHorizModule,
      minimum height=16\TPVertModule] {};%
  \end{textblock}%
}\makeatother
%    \end{macrocode}
% \end{macro}

% \begin{macro}{\huawei@logo}
% Then, we define |\huawei@logo|:
%    \begin{macrocode}
\makeatletter\newcommand\huawei@logo{
\def\BLACK{\ifhuawei@dark white\else black\fi}
\def\BLANK{\ifhuawei@dark black\else white\fi}
\def\SCALE{\ifhuawei@slides 1.8\else 1\fi}
\begin{tikzpicture}[y=0.80pt, x=0.80pt, yscale=-\SCALE,
xscale=\SCALE, inner sep=0pt, outer sep=0pt]
\begin{scope}[even odd rule,line width=0.800pt]
\begin{scope}[shift={(0,-0.00024)}]
\path[fill=\BLACK] (10.9375,30.2240) -- (10.9375,33.6097)
.. controls (10.9375,34.5713) and (10.4603,35.0845) ..
   (9.5932,35.0845) .. controls (8.7212,35.0845) and
   (8.2411,34.5565)
.. (8.2411,33.5686) -- (8.2411,30.2278) -- (7.0415,30.2278)
-- (7.0415,33.6097) .. controls (7.0415,35.2737) and
   (7.9658,36.2272) .. (9.5774,36.2272) .. controls
   (11.2041,36.2272) and (12.1371,35.2554) .. (12.1371,33.5609) --
   (12.1371,30.2240)
-- (10.9375,30.2240) -- cycle;
\path[fill=\BLACK] (15.3511,30.2240) -- (12.7456,36.1351) --
 (13.9702,36.1351) -- (14.4731,34.9903) -- (14.5091,34.9045) --
 (17.2158,34.9045) -- (17.7467,36.1351) -- (19.0045,36.1351) --
 (16.4233,30.2590) -- (16.4001,30.2240) -- cycle;
\path[fill=\BLACK] (22.4840,30.2240) -- (21.1414,34.2912) --
 (19.8344,30.2271) -- (18.5578,30.2271) -- (20.6186,36.1388) --
 (21.6120,36.1388) -- (22.9573,32.2553) -- (24.3016,36.1388) --
 (25.3034,36.1388) -- (27.3592,30.2271) -- (26.1152,30.2271) --
 (24.8055,34.2912) -- (23.4626,30.2240) -- cycle;
\path[fill=\BLACK] (34.2236,30.2240) -- (34.2236,36.1300) --
 (35.4074,36.1300) -- (35.4074,30.2240) -- cycle;
\path[fill=\BLACK] (0.2686,30.2244) -- (0.2686,36.1384) --
 (1.4686,36.1384) -- (1.4686,33.7365) -- (4.1780,33.7365) --
 (4.1780,36.1384) -- (5.3783,36.1384) -- (5.3783,30.2244) --
 (4.1780,30.2244) -- (4.1780,32.6102) -- (1.4686,32.6102) --
 (1.4686,30.2244) -- cycle;
\path[fill=\BLACK] (28.3267,30.2284) -- (28.3267,36.1344) --
 (32.7928,36.1344) -- (32.7928,35.0575) -- (29.5105,35.0575) --
 (29.5105,33.5931) -- (31.6931,33.5931) -- (31.6931,32.5160) --
 (29.5105,32.5160) -- (29.5105,31.3052) -- (32.6785,31.3052) --
 (32.6785,30.2284) -- cycle;
\path[fill=\BLANK] (15.8594,31.7207) -- (16.7149,33.7008) --
 (16.7119,33.7008) -- (16.7701,33.8374) -- (14.9552,33.8374) --
 (15.0127,33.7008) -- (15.0117,33.7008) -- cycle;
\path[fill=red] (15.2113,0.0001) .. controls (14.7353,0.0422) and
 (13.4491,0.3349) .. (13.4491,0.3349) -- (13.4486,0.3349) .. controls
 (10.5495,1.0843) and (9.8643,3.7151) .. (9.8643,3.7151) .. controls
 (9.7227,4.1575) and (9.6584,4.6109) .. (9.6395,5.0380) --
 (9.6395,5.6194) .. controls (9.6780,6.5176) and (9.8777,7.1883) ..
 (9.8777,7.1883) .. controls (10.8461,11.4823) and
 (15.6075,18.5379) ..(16.6308,20.0200) .. controls
 (16.7034,20.0920) and(16.7615,20.0660) ..
 (16.7615,20.0660) .. controls(16.8720,20.0355) and
 (16.8633,19.9296) .. (16.8633,19.9296) --
 (16.8654,19.9301) .. controls (18.4416,4.1760) and
 (15.2113,0.0001) .. (15.2113,0.0001) -- (15.2113,0.0001) -- cycle;
\path[fill=red] (20.4192,0.0000) .. controls (20.4192,0.0000) and
 (17.1748,4.1782) .. (18.7521,19.9415) --
 (18.7542,19.9415) .. controls (18.7664,20.0416) and
 (18.8374,20.0624) ..(18.8374,20.0624) .. controls
 (18.9430,20.1033) and(18.9971,20.0019) .. (18.9971,20.0019) --
 (18.9976,20.0029) .. controls (20.0472,18.4829) and
 (24.7814,11.4657) ..(25.7455,7.1882) .. controls(25.7455,7.1882) and
 (26.2683,5.1179) .. (25.7636,3.7150) .. controls
 (25.7636,3.7150) and (25.0470,1.0434) ..(22.1442,0.3375) .. controls
 (22.1442,0.3375) and (21.3077,0.1251) .. (20.4193,0.0000) --
 (20.4192,0.0000) -- cycle;
\path[fill=red] (5.4542,4.7294) .. controls (5.4542,4.7294) and
 (2.6920,7.3513) .. (2.5547,10.1307) -- (2.5558,10.1307) --
 (2.5558,10.5519) .. controls (2.5579,10.5836) and
 (2.5595,10.6157) .. (2.5615,10.6480) .. controls
 (2.6804,12.8893) and (4.3665,14.2157) ..
 (4.3665,14.2157) .. controls (7.0831,16.8634) and
 (13.6611,20.2062) .. (15.1896,20.9647) .. controls
 (15.2110,20.9727) and (15.2900,20.9985) ..
 (15.3364,20.9399) .. controls (15.3364,20.9399) and
 (15.3605,20.9216) .. (15.3715,20.8887) --
 (15.3715,20.8189) .. controls (15.3704,20.8149) and
 (15.3683,20.8109) .. (15.3663,20.8065) --
 (15.3669,20.8065) .. controls (11.1809,11.6620) and
 (5.4543,4.7294) .. (5.4543,4.7294) -- (5.4542,4.7294) -- cycle;
\path[fill=red] (30.1695,4.7294) .. controls (30.1695,4.7294) and
 (24.4602,11.6403) .. (20.2761,20.7662) --
 (20.2777,20.7657) .. controls (20.2777,20.7657) and
 (20.2283,20.8715) .. (20.3082,20.9399) .. controls
 (20.3082,20.9399) and (20.3320,20.9579) .. (20.3650,20.9652) --
 (20.4229,20.9652) .. controls (20.4333,20.9622) and
 (20.4444,20.9582) .. (20.4554,20.9507) --
 (20.4554,20.9517) .. controls (22.0255,20.1724) and
 (28.5525,16.8516) .. (31.2563,14.2162) .. controls
 (31.2563,14.2162) and (32.9686,12.8412) ..
 (33.0583,10.6334) .. controls (33.2560,7.5672) and
 (30.1696,4.7294) .. (30.1696,4.7294) -- (30.1695,4.7294) -- cycle;
\path[fill=red] (35.3082,13.8080) .. controls (35.3082,13.8080) and
 (26.0003,18.8031) .. (21.1876,22.0494) -- (21.1882,22.0499) --
 (21.1893,22.0509) .. controls (21.1893,22.0509) and
 (21.1019,22.1081) .. (21.1324,22.2106) .. controls
 (21.1324,22.2106) and (21.1781,22.2933) .. (21.2446,22.2933) --
 (21.2446,22.2938) .. controls (22.9684,22.2968) and
 (29.4914,22.3041) .. (29.6549,22.2736) .. controls
 (29.6549,22.2736) and (30.4983,22.2399) ..
 (31.5406,21.8395) .. controls (31.5406,21.8395) and
 (33.8608,21.1019) .. (35.0659,18.4677) .. controls
 (35.0659,18.4677) and (35.6166,17.3662) .. (35.6246,15.8187) --
 (35.6246,15.7531) .. controls (35.6206,15.1563) and
 (35.5365,14.4957) .. (35.3083,13.8080) --
 (35.3082,13.8080) -- cycle;
\path[fill=red] (0.3129,13.8313) .. controls (-0.5339,16.4531) and
 (0.6062,18.5656) .. (0.6080,18.5690) .. controls
 (1.7962,21.0795) and (4.0641,21.8401) ..
 (4.0641,21.8401) .. controls (5.1112,22.2703) and
 (6.1586,22.3000) .. (6.1586,22.3000) .. controls
 (6.3221,22.3300) and (12.6713,22.3030) ..
 (14.3725,22.2950) .. controls (14.4445,22.2945) and
 (14.4842,22.2221) .. (14.4842,22.2221) .. controls
 (14.4892,22.2142) and (14.4936,22.2062) .. (14.4955,22.1989) --
 (14.4955,22.1333) .. controls (14.4814,22.0893) and
 (14.4439,22.0584) .. (14.4439,22.0584) --
 (14.4450,22.0578) .. controls (9.6347,18.8126) and
 (0.3130,13.8315) .. (0.3130,13.8315) -- (0.3129,13.8313) -- cycle;
\path[fill=red] (14.1782,23.1010) -- (3.2285,23.4850) .. controls
 (4.4160,25.6018) and (6.4158,27.2469) ..
 (8.4990,26.7416) .. controls (9.9366,26.3822) and
 (13.1935,24.1105) .. (14.2687,23.3429) --
 (14.2651,23.3399) .. controls (14.3487,23.2648) and
 (14.3188,23.2045) .. (14.3188,23.2045) .. controls
 (14.2913,23.1062) and (14.1783,23.1063) .. (14.1783,23.1063) --
 (14.1782,23.1010) -- cycle;
\path[fill=red] (21.4305,23.1090) -- (21.4294,23.1130) .. controls
 (21.4294,23.1130) and (21.3333,23.1252) ..
 (21.3064,23.1972) .. controls (21.3064,23.1972) and
 (21.2831,23.2956) .. (21.3478,23.3445) --
 (21.3467,23.3455) .. controls (22.3950,24.0973) and
 (25.5729,26.3191) .. (27.1039,26.7505) .. controls
 (27.1039,26.7505) and (27.3359,26.8294) .. (27.7246,26.8435) --
 (27.9581,26.8435) .. controls (28.9839,26.8075) and
 (30.7708,26.2805) .. (32.3946,23.4907) --
 (21.4305,23.1090) -- cycle;
\end{scope}%
\end{scope}%
\end{tikzpicture}%
}\makeatother
%    \end{macrocode}
% \end{macro}

% Then, we configure headers using |fancyhdr|:
%    \begin{macrocode}
\RequirePackage{fancyhdr}
\pagestyle{fancy}
\renewcommand{\headrulewidth}{0pt}
\fancyhf{}
\makeatletter\fancyfoot[L]{
  \huawei@bar
  \ifhuawei@authordraft%
    \begin{textblock}{14}[0.5,0.5](8,8)%
      \tikz \node[minimum width=14\TPHorizModule] {%
        \fontsize{64}{64}\selectfont\sffamily\scshape%
        \color{gray!20}\rotatebox{30}{it is a draft}
      };%
    \end{textblock}%
  \else\fi%
}\makeatother
\makeatletter\fancyhead[L]{
  \ifnum\value{page}=1\else%
    \ifhuawei@nobrand\else%
      \begin{textblock}{8}[0,0](1.2,0.2)%
        \huawei@logo%
      \end{textblock}%
    \fi%
  \fi%
}\makeatother
\makeatletter\fancyhead[R]{
  \begin{textblock}{8}[1,0](15.8,0.2)%
    \raggedleft\huawei@header{%
      \ifhuawei@nosecurity\else%
        \thesecurity
      \fi%
    }%
  \end{textblock}%
}\makeatother
\makeatletter\fancyfoot[R]{
  \begin{textblock}{8}[0,1](1.2,15.8)%
    \ifnum\value{page}=1\else%
      \huawei@header{\raggedright%
        \ifhuawei@anonymous\else%
          \theauthor%
          \ifhuawei@nosecurity\else%
            \ifx\theid\empty\else, \theid\fi
          \fi%
          \ifhuawei@nobrand\else%
            \newline
          \fi%
        \fi%
        \ifhuawei@nobrand\else%
          \thecompany{}
        \fi%
      }%
    \fi%
  \end{textblock}%
  \begin{textblock}{8}[1,1](15.8,15.8)%
    \raggedleft\huawei@header{%
      \ifnum\value{page}=1\else%
        \ifhuawei@nopaging\else%
          Page \#\thepage{} of \pageref*{LastPage}%
          \ifhuawei@nodate\else%
            \\
          \fi%
        \fi%
      \fi%
      \ifhuawei@nodate\else%
        \today{} \settimeformat{daytime}\currenttime{}%
      \fi%
    }%
  \end{textblock}%
}\makeatother
%    \end{macrocode}

% \begin{macro}{abstract}
% Then, we redefine |abstract| environment:
%    \begin{macrocode}
\RequirePackage{changepage}
\renewenvironment{abstract}
  {\begin{adjustwidth}{0pt}{1in}{\scshape Abstract:}\newline\small}
  {\end{adjustwidth}}
%    \end{macrocode}
% \end{macro}

% \begin{macro}{\maketitle}
% Then, we redefine |\maketitle|:
%    \begin{macrocode}
\makeatletter\renewcommand\maketitle{%
  \vspace*{18pt}
  {\bfseries{\Huge\thetitle}}
  \ifx\thesubtitle\empty\else%
    \\[6pt]
    {\color{gray!50!black}\large\thesubtitle}
  \fi
  \\[18pt]
  \ifhuawei@nobrand\else%
    \ifhuawei@anonymous\else%
      \ifx\thecompany\empty\else%
        \thecompany\newline
      \fi
    \fi
  \fi
  \ifx\theauthor\empty\else%
    {\scshape\ifhuawei@anonymous
      Anonymous Authors
    \else%
      \theauthor
    \fi}%
  \fi%
  \ifhuawei@anonymous\else%
    \ifx\theauthor\empty\else
      \ifhuawei@nobrand\else%
        \ifx\theid\empty\else
          $\;$/ {\theid}
        \fi
      \fi
    \fi
  \fi
  \vspace{2em}
}\makeatother
%    \end{macrocode}
% \end{macro}

% \begin{macro}{\PrintCrumb}
% Then, we define |\PrintCrumb|:
%    \begin{macrocode}
\newcommand\PrintCrumb[2]{%
  \begin{minipage}{\columnwidth}%
    \raggedright\textsc{#1}:\\#2%
  \end{minipage}\vspace{4pt}%
}
%    \end{macrocode}
% \end{macro}

% \begin{macro}{\PrintThankYouPage}
% Then, we define |\PrintThankYouPage|:
%    \begin{macrocode}
\newcommand\PrintThankYouPage{
  \newpage
  \vspace*{\fill}
  \begin{center}
    \normalsize
    {\Huge\color{red}\textbf{Thank you!}}
  \end{center}
  \vspace*{\fill}
}
%    \end{macrocode}
% \end{macro}

% \begin{macro}{\PrintDisclaimer}
% Then, we define |\PrintDisclaimer|:
%    \begin{macrocode}
\makeatletter\newcommand\PrintDisclaimer{%
  \justify\vspace*{\fill}%
  \begingroup%
  \setstretch{0.55}%
  \sffamily\scriptsize\color{gray!50!black}%
  \textbf{Disclaimer}: The opinions expressed in
  this document are in good faith and
  while every care has been taken in preparing it,
  \ifhuawei@nobrand%
    the author%
  \else%
    \thecompany{}%
  \fi
  makes no representations and gives no warranties
  of whatever nature in respect of these documents,
  including but not limited to the accuracy or completeness
  of any information, facts and/or opinions contained therein.
  \ifhuawei@nobrand%
    The author%
  \else%
    \thecompany{}%
  \fi,
  its subsidiaries, the directors, employees and agents
  cannot be held liable for the use of and reliance of
  the opinions, estimates, forecasts and findings in
  these documents.
  \par
  \endgroup%
}\makeatother
%    \end{macrocode}
% \end{macro}

% \begin{macro}{\PrintFirstPage}
% Then, we define |\PrintFirstPage|:
%    \begin{macrocode}
\makeatletter\newcommand\PrintFirstPage[1]{
  \ifhuawei@landscape\else
    \PackageError{huawei}{It's allowed to use
    PrintFirstPage only in landscape mode}{Read huawei.pdf
    for more information}
  \fi
  \huawei@bar
  \def\param{#1}%
  \ifx\param\empty\else
    \includegraphics[height=2in]{#1}
    \newline
  \fi
  \vspace*{0.5in}
  \maketitle
  \ifhuawei@nocover\else
    \begin{textblock}{14}[1,1](14,14)
      \raggedleft\includegraphics[height=3.6in]{huawei-cover-picture.pdf}
    \end{textblock}
  \fi
}\makeatother
%    \end{macrocode}
% \end{macro}

% \begin{macro}{\PrintLastPage}
% Then, we define |\PrintLastPage|:
%    \begin{macrocode}
\makeatletter\newcommand\PrintLastPage{
  \ifhuawei@landscape\else
    \PackageError{huawei}{It's allowed to use
    PrintLastPage only in landscape mode}{Read huawei.pdf
    for more information}
  \fi
  \newpage
  \vspace*{1in}
  \begin{center}
    \begin{minipage}{0.6\columnwidth}\raggedright
      \normalsize
      \setlength{\parskip}{6pt}
      {\Huge\color{red}\textbf{Thank you!}}

      \vspace{0.5in}

      Bring digital to every person, home and organization
      \newline
      for a fully connected, intelligent world.

      \textbf{%
        Copyright \copyright{} \the\year{} \thecompany{}
        \newline
        All Rights Reserved.%
      }

      The information in this document may contain predictive
      statements including, without limitation, statements regarding
      the future financial and operating results, future product
      portfolio, new technology, etc. There are a number of factors that
      could cause actual results and developments to differ materially
      from those expressed or implied in the predictive statements.
      Therefore, such information is provided for reference purpose
      only and constitutes neither an offer nor an acceptance.
      \ifhuawei@nobrand%
        The author%
      \else%
        \thecompany{}%
      \fi
      may change the information at any time without notice.
    \end{minipage}
  \end{center}%
}\makeatother
%    \end{macrocode}
% \end{macro}

% \begin{macro}{\PrintBibliography}
% Then, we define |\PrintBibliography|:
%    \begin{macrocode}
\makeatletter\newcommand\PrintBibliography{
  \setlength\bibitemsep{3pt}
  \AtNextBibliography{\small}
  \newpage
  \raggedright
  \ifhuawei@landscape
    \begin{multicols}{3}\setstretch{0.95}\printbibliography\end{multicols}
  \else
    \begin{multicols}{2}\setstretch{0.95}\printbibliography\end{multicols}
  \fi
}\makeatother
%    \end{macrocode}
% \end{macro}

% Then, we set the background color of the document:
%    \begin{macrocode}
\makeatletter\ifhuawei@dark
  \RequirePackage{pagecolor}
  \pagecolor{black}
  \color{white}
\fi
%    \end{macrocode}

% Then, we configure layout:
%    \begin{macrocode}
\AtBeginDocument{%
\raggedbottom%
\raggedcolumns%
\setlength\headheight{32pt}%
\setlength\footskip{32pt}%
\setlength\topskip{0mm}%
\setlength\parindent{0pt}%
\setlength\parskip{6pt}%
\setlength\columnsep{32pt}%
\def\arraystretch{1.5}%
}
%    \end{macrocode}

% Then, we renew a few commands:
%    \begin{macrocode}
\renewcommand\title[1]{\renewcommand\thetitle{#1}}
\newcommand\thetitle{\textbackslash{}thetitle}
\newcommand*\thecompany{Huawei Technologies Co., Ltd.}
\newcommand\thesubtitle{}
\renewcommand\author[1]{\renewcommand\theauthor{#1}}
\newcommand\theauthor{\textbackslash{}theauthor}
\newcommand*\thesecurity{Confidential}
\newcommand*\theid{}
\ifcsname nospell\endcsname\else\newcommand\nospell[1]{#1}\fi
%    \end{macrocode}

% \Finale

%\clearpage
%
%\PrintChanges
%\clearpage
%\PrintIndex
